\PassOptionsToPackage{table}{xcolor}
\documentclass[10pt,aspectratio=169]{beamer}
\newcommand{\LectureNumber}{22}
\input{course-theme.tex}
\title{Two-dimensional arrays}
\subtitle{CSC103 Programming Fundamentals / Lecture 22}
\author{Dr. Muhammad Farhan}
\institute{Associate Professor of Computer Science\\Department of Computer Science\\COMSATS University Islamabad\\Sahiwal Campus}
\date{Fall 2026}
\begin{document}
\begin{frame}\titlepage\end{frame}

\begin{frame}[fragile,t]{The problem for this lecture}
\begin{columns}[T,onlytextwidth]\begin{column}{0.60\textwidth}
Compute the total of every element in \{\{2,4\},\{6,8\}\}. Then compute its main diagonal sum.
\end{column}\begin{column}{0.35\textwidth}\small
Traverse every element with nested loops to accumulate the grand total.
\end{column}\end{columns}
\note{Introduce the target problem before explaining the tools. Revisit it in the guided solution later.\par Syllabus reading: Deitel Ch 7. CLO-2.}
\end{frame}

\begin{frame}[fragile,t]{Learning objectives}
\begin{columns}[T,onlytextwidth]\begin{column}{0.60\textwidth}
\begin{itemize}
\item Create and access a table of values
\item Traverse rows and columns
\item Compute row and column totals
\end{itemize}
\end{column}\begin{column}{0.35\textwidth}\small
CLO-2

Arrays of arrays\par Nested traversal\par Row totals\par Column totals
\end{column}\end{columns}
\note{Explain how each objective supports the opening problem. The final questions check these objectives.\par Syllabus reading: Deitel Ch 7. CLO-2.}
\end{frame}

\begin{frame}[fragile,t]{Arrays of arrays}
\begin{itemize}
\item A Java two-dimensional array is an array of row references.
\item data[r][c] selects column c in row r.
\item A rectangular example has equal row lengths.
\end{itemize}
\vspace{1mm}\begin{block}{Example}
\begin{lstlisting}
int[][] marks = {
  {60, 70, 80},
  {50, 65, 75}
};
\end{lstlisting}
\end{block}
\note{The outer array has length 2. Each row in this example has length 3. Java does not require every row to share that length.\par Syllabus reading: Deitel Ch 7. CLO-2.}
\end{frame}

\begin{frame}[fragile,t]{Nested traversal}
\begin{itemize}
\item Use the outer length for row count.
\item Use each row length for its column count.
\item Keep the row and column indexes distinct.
\end{itemize}
\vspace{1mm}\begin{block}{Example}
\begin{lstlisting}
for (int r=0; r<data.length; r++) {
  for (int c=0; c<data[r].length; c++) {
    System.out.println(data[r][c]);
  }
}
\end{lstlisting}
\end{block}
\note{This traversal accommodates non-null ragged rows. A null row needs separate handling.\par Syllabus reading: Deitel Ch 7. CLO-2.}
\end{frame}

\begin{frame}[fragile,t]{Row totals}
\begin{itemize}
\item Reset a row accumulator at the start of each row.
\item The inner loop adds elements from that row.
\item A row mean divides by the row length only when nonempty.
\end{itemize}
\vspace{1mm}\begin{block}{Example}
\begin{lstlisting}
for (int[] row : data) {
  int total = 0;
  for (int value : row) total += value;
  System.out.println(total);
}
\end{lstlisting}
\end{block}
\note{Explain that placing total outside this outer loop would produce cumulative totals rather than independent row totals.\par Syllabus reading: Deitel Ch 7. CLO-2.}
\end{frame}

\begin{frame}[fragile,t]{Column totals}
\begin{itemize}
\item Column traversal reverses the roles of the indexes.
\item A rectangular-matrix contract simplifies column totals.
\item Check shape assumptions before using data[0].
\end{itemize}
\vspace{1mm}\begin{block}{Example}
\begin{lstlisting}
For {{1,2},{3,4}}:
Column 0 total: 1 + 3 = 4
Column 1 total: 2 + 4 = 6
\end{lstlisting}
\end{block}
\note{For an empty outer array data[0] is invalid. A ragged array can lack a requested column in some rows.\par Syllabus reading: Deitel Ch 7. CLO-2.}
\end{frame}

\begin{frame}[fragile,t]{Different traversals answer different questions}
\begin{itemize}
\item A row traversal keeps the row fixed while the column changes.
\item A column traversal keeps the column fixed while the row changes.
\item A main-diagonal traversal selects entries whose row and column indexes match.
\end{itemize}
\vspace{1mm}\begin{block}{Example}
\begin{lstlisting}
For {{2,4},{6,8}}:
Row 0: 2,4
Column 0: 2,6
Main diagonal: 2,8
\end{lstlisting}
\end{block}
\note{Explain the mechanism using the displayed example. A row traversal keeps the row fixed while the column changes. A column traversal keeps the column fixed while the row changes. A main-diagonal traversal selects entries whose row and column indexes match.\par Syllabus reading: Deitel Ch 7. CLO-2.}
\end{frame}

\begin{frame}[fragile,t]{Accumulator placement changes the meaning}
\begin{itemize}
\item A grand total starts before both loops.
\item A row total resets inside the outer loop.
\item A diagonal sum needs one selected element per suitable row.
\end{itemize}
\vspace{1mm}\begin{block}{Example}
\begin{lstlisting}
Grand total: 2 + 4 + 6 + 8 = 20
Diagonal total: 2 + 8 = 10
\end{lstlisting}
\end{block}
\note{Explain the mechanism using the displayed example. A grand total starts before both loops. A row total resets inside the outer loop. A diagonal sum needs one selected element per suitable row.\par Syllabus reading: Deitel Ch 7. CLO-2.}
\end{frame}


\begin{frame}[fragile,t]{Rows and columns identify one element}
\begin{center}\begin{tikzpicture}
\node[box,text width=13mm,minimum height=11mm] (n0) at (0.0,0.0) {2};
\node[box,text width=13mm,minimum height=11mm] (n1) at (1.8,0.0) {4};
\node[box,text width=13mm,minimum height=11mm] (n2) at (3.6,0.0) {6};
\node[box,text width=13mm,minimum height=11mm] (n3) at (0.0,-1.2) {1};
\node[box,text width=13mm,minimum height=11mm] (n4) at (1.8,-1.2) {3};
\node[box,text width=13mm,minimum height=11mm] (n5) at (3.6,-1.2) {5};
\node[font=\small,text=accent] at (0.0,.85) {column 0};
\node[font=\small,text=accent] at (1.8,.85) {column 1};
\node[font=\small,text=accent] at (3.6,.85) {column 2};
\node[font=\small,text=accent] at (-1.7,0.0) {row 0};
\node[font=\small,text=accent] at (-1.7,-1.2) {row 1};
\end{tikzpicture}\end{center}
\begin{block}{What to notice}For this 2 by 3 array, data[1][2] is 5: row 1, column 2.\end{block}
\note{Trace each labelled connection. For this 2 by 3 array, data[1][2] is 5: row 1, column 2.}
\end{frame}
\begin{frame}[fragile,t]{Key distinctions}
\small\renewcommand{\arraystretch}{1.5}
\rowcolors{2}{pale}{white}
\begin{tabularx}{\textwidth}{>{\raggedright\arraybackslash}X>{\raggedright\arraybackslash}X>{\raggedright\arraybackslash}X>{\raggedright\arraybackslash}X}
\rowcolor{navy}
\color{white}\bfseries Row / column & \color{white}\bfseries Column 0 & \color{white}\bfseries Column 1 & \color{white}\bfseries Row total\\
Row 0 & 2 & 4 & 6\\
Row 1 & 6 & 8 & 14\\
Column total & 8 & 12 & 20 overall\\
\end{tabularx}
\vspace{3mm}\footnotesize These distinctions determine which operation or control structure fits the problem.
\note{\par Syllabus reading: Deitel Ch 7. CLO-2.}
\end{frame}

\begin{frame}[fragile,t]{First example: execution}
\begin{lstlisting}
int[][] data = {{1, 2, 3}, {4, 5, 6}};
for (int[] row : data) {
  int total = 0;
  for (int value : row) total += value;
  System.out.println(total);
}
\end{lstlisting}
\note{This is the main body from Java\_Examples/Lecture22.java. The source file includes the complete class. Predict the result before the next slide.\par Syllabus reading: Deitel Ch 7. CLO-2.}
\end{frame}

\begin{frame}[fragile,t]{First example: result}
\begin{columns}[T,onlytextwidth]\begin{column}{0.60\textwidth}
6\par 15\par A fresh total starts at zero for each row.
\end{column}\begin{column}{0.35\textwidth}\small
Row totals

Column totals
\end{column}\end{columns}
\note{Expected output and interpretation: 6
15
A fresh total starts at zero for each row.\par Syllabus reading: Deitel Ch 7. CLO-2.}
\end{frame}

\begin{frame}[fragile,t]{Guided practice}
\begin{columns}[T,onlytextwidth]\begin{column}{0.60\textwidth}
Compute the total of every element in \{\{2,4\},\{6,8\}\}. Then compute its main diagonal sum.
\end{column}\begin{column}{0.35\textwidth}\small
Attempt the solution before continuing.

Use the definitions and examples from this lecture.
\end{column}\end{columns}
\note{Give students 8 minutes to attempt the problem. The following slides provide a complete solution. Original solution guidance: All-element sum: 20. Main diagonal: data[0][0]+data[1][1]=10. For general diagonal traversal use the declared matrix shape contract.\par Syllabus reading: Deitel Ch 7. CLO-2.}
\end{frame}

\begin{frame}[fragile,t]{Solution strategy}
\begin{columns}[T,onlytextwidth]\begin{column}{0.60\textwidth}
\begin{itemize}
\item Traverse every element with nested loops to accumulate the grand total.
\item For this square matrix, add data[i][i] in a separate loop.
\item Report each result with a distinct label.
\end{itemize}
\end{column}\begin{column}{0.35\textwidth}\small
The next program uses fixed inputs so its result can be reproduced.
\end{column}\end{columns}
\note{Discuss why each step is needed. State the input assumptions before executing the program.\par Syllabus reading: Deitel Ch 7. CLO-2.}
\end{frame}

\begin{frame}[fragile,t]{Complete solution (1/2)}
\begin{lstlisting}
public class Solution22 {
  public static void main(String[] args) throws Exception {
    int[][] data = {{2, 4}, {6, 8}};
    int total = 0;
    for (int[] row : data) {
      for (int value : row) total += value;
    }
    int diagonal = 0;
\end{lstlisting}
\note{Complete runnable file: Java\_Examples/Solution22.java. Inputs are fixed in the program.  \par Syllabus reading: Deitel Ch 7. CLO-2.}
\end{frame}

\begin{frame}[fragile,t]{Complete solution (2/2)}
\begin{lstlisting}
    for (int i = 0; i < data.length; i++) {
      diagonal += data[i][i];
    }
    System.out.println("Total: " + total);
    System.out.println("Diagonal: " + diagonal);
  }
}
\end{lstlisting}
\note{Complete runnable file: Java\_Examples/Solution22.java. Inputs are fixed in the program.  \par Syllabus reading: Deitel Ch 7. CLO-2.}
\end{frame}

\begin{frame}[fragile,t]{Execution trace}
\small\renewcommand{\arraystretch}{1.5}
\rowcolors{2}{pale}{white}
\begin{tabularx}{\textwidth}{>{\raggedright\arraybackslash}X>{\raggedright\arraybackslash}X>{\raggedright\arraybackslash}X}
\rowcolor{navy}
\color{white}\bfseries Selected element & \color{white}\bfseries Value & \color{white}\bfseries Running grand total\\
data[0][0] & 2 & 2\\
data[0][1] & 4 & 6\\
data[1][0] & 6 & 12\\
data[1][1] & 8 & 20\\
\end{tabularx}
\vspace{3mm}\footnotesize Follow the changing state in execution order.
\note{Manually derived trace for the complete solution. Traverse every element with nested loops to accumulate the grand total. For this square matrix, add data[i][i] in a separate loop. Report each result with a distinct label.\par Syllabus reading: Deitel Ch 7. CLO-2.}
\end{frame}

\begin{frame}[fragile,t]{Solution result}
\begin{columns}[T,onlytextwidth]\begin{column}{0.60\textwidth}
Total: 20\par Diagonal: 10
\end{column}\begin{column}{0.35\textwidth}\small
Why this result follows

Report each result with a distinct label.
\end{column}\end{columns}
\note{Expected result: Total: 20
Diagonal: 10\par Syllabus reading: Deitel Ch 7. CLO-2.}
\end{frame}

\begin{frame}[fragile,t]{A common incorrect approach}
for(int c=0; c\textless{}data.length; c++) \{\par   total += data[r][c];\par \}
\vspace{1mm}\begin{block}{Example}
\begin{lstlisting}
The bound uses the number of rows. Use data[r].length for columns in row r.
\end{lstlisting}
\end{block}
\note{Ask students to predict the failure before discussing the explanation. The right side gives the correction.\par Syllabus reading: Deitel Ch 7. CLO-2.}
\end{frame}

\begin{frame}[fragile,t]{Boundary and failure checks}
\begin{columns}[T,onlytextwidth]\begin{column}{0.60\textwidth}
Check the stated input assumptions.

Represent marks for three students in two quizzes. Print each student total and each quiz average.
\end{column}\begin{column}{0.35\textwidth}\small
All-element sum: 20. Main diagonal: data[0][0]+data[1][1]=10. For general diagonal traversal use the declared matrix shape contract.
\end{column}\end{columns}
\note{Use the specification to derive each expected result. Exercise solution: All-element sum: 20. Main diagonal: data[0][0]+data[1][1]=10. For general diagonal traversal use the declared matrix shape contract.\par Syllabus reading: Deitel Ch 7. CLO-2.}
\end{frame}

\begin{frame}[fragile,t]{Check your understanding}
\begin{columns}[T,onlytextwidth]\begin{column}{0.60\textwidth}
What does data.length measure? Where should a row-total accumulator be reset?
\end{column}\begin{column}{0.35\textwidth}\small
Write a prediction and explain the rule that supports it.
\end{column}\end{columns}
\note{Answer: The number of rows in the outer array. Reset it inside the outer row loop before the inner traversal.\par Syllabus reading: Deitel Ch 7. CLO-2.}
\end{frame}

\begin{frame}[fragile,t]{Answer and explanation}
\begin{columns}[T,onlytextwidth]\begin{column}{0.60\textwidth}
The number of rows in the outer array. Reset it inside the outer row loop before the inner traversal.
\end{column}\begin{column}{0.35\textwidth}\small
The bound uses the number of rows. Use data[r].length for columns in row r.
\end{column}\end{columns}
\note{Reconnect the answer to the relevant definition rather than treating it as a fact to memorise.\par Syllabus reading: Deitel Ch 7. CLO-2.}
\end{frame}


\begin{frame}[fragile,t]{Compare cases and predict outcomes}
\small\renewcommand{\arraystretch}{1.5}\rowcolors{2}{pale}{white}
\begin{tabularx}{\textwidth}{>{\raggedright\arraybackslash}X>{\raggedright\arraybackslash}X>{\raggedright\arraybackslash}X}\rowcolor{navy}
\color{white}\bfseries Aggregate & \color{white}\bfseries Elements & \color{white}\bfseries Result\\
Row 0 & 2 + 4 + 6 & 12\\
Column 1 & 4 + 3 & 7\\
All elements & 2 + 4 + 6 + 1 + 3 + 5 & 21\\
\end{tabularx}\par\vspace{3mm}\begin{exampleblock}{Discuss}Explain each expected result before running the example. Which case would reveal a wrong assumption?\end{exampleblock}
\note{Read the first two columns and invite predictions before discussing the outcome.}
\end{frame}

\begin{frame}[fragile,t]{Apply the idea}
\begin{alertblock}{Independent practice}For \{\{2,4,6\},\{1,3,5\}\}, compute column 1 using a loop over rows. State the shape assumption.\end{alertblock}\vspace{3mm}\begin{enumerate}\item Work through the input by hand.\item Write or correct the program.\item Justify the expected result.\end{enumerate}\note{Allow 3–5 minutes, then compare answers in pairs. For \{\{2,4,6\},\{1,3,5\}\}, compute column 1 using a loop over rows. State the shape assumption.}
\end{frame}

\begin{frame}[fragile,t]{Solution and reasoning}
\begin{lstlisting}
int[][] data = {{2, 4, 6}, {1, 3, 5}};
int total = 0;
for (int row = 0; row < data.length; row++) {
  total += data[row][1];
}
System.out.println(total);
\end{lstlisting}\vspace{1mm}\small\begin{itemize}
\item Column 1 contains 4 and 3, so the output is 7.
\item Every visited row must be non-null and contain index 1.
\item For arbitrary ragged data, check each row before accessing that column.
\end{itemize}\note{Answer: Column 1 contains 4 and 3, so the output is 7. Every visited row must be non-null and contain index 1. For arbitrary ragged data, check each row before accessing that column.}
\end{frame}

\begin{frame}[fragile,t]{Grading a multiple{-}choice test}
\begin{itemize}
\item Rows represent students and columns represent questions.
\item Compare each answer with the key at the same column index.
\item Reset the score for each student, so scores do not accumulate across rows.
\end{itemize}
\vspace{3mm}\footnotesize\textcolor{accent}{Source: supplied ISB campus slides. See speaker notes and ISB\_Source\_Map.md.}
\note{Adapted from the supplied ISB campus folder: Multi{-}dimensional arrays.pptx, slides 71, 72. Uses the first three of the ten supplied students; retains the original five{-}question key and answers. Replaced typographic quotes with Java quotes.}
\end{frame}

\begin{frame}[fragile,t]{Worked problem: Grading a multiple{-}choice test}
\begin{block}{Task}Use the first three source students and answer key \{D,B,D,C,C\}. Compute each score out of five.\end{block}
\small Input: Values are fixed in the program for a reproducible demonstration.\par\vspace{3mm}\begin{exampleblock}{Before running}Predict the output and identify the variables that change.\end{exampleblock}
\note{Adapted from the supplied ISB campus folder: Multi{-}dimensional arrays.pptx, slides 71, 72. Uses the first three of the ten supplied students; retains the original five{-}question key and answers. Replaced typographic quotes with Java quotes.}
\end{frame}

\begin{frame}[fragile,t]{Complete ISB example (1/2)}
\begin{lstlisting}
public class IsbLecture22 {
  public static void main(String[] args) throws Exception {
    char[][] answers = {
      {'A','B','A','C','C'},
      {'D','B','A','B','C'},
      {'E','D','D','A','C'}
    };
    char[] key = {'D','B','D','C','C'};
    for (int student = 0; student < answers.length; student++) {
\end{lstlisting}
\note{Adapted from the supplied ISB campus folder: Multi{-}dimensional arrays.pptx, slides 71, 72. Uses the first three of the ten supplied students; retains the original five{-}question key and answers. Replaced typographic quotes with Java quotes. Complete file: ISB\_Java\_Examples/IsbLecture22.java.}
\end{frame}

\begin{frame}[fragile,t]{Complete ISB example (2/2)}
\begin{lstlisting}
      int score = 0;
      for (int question = 0; question < key.length; question++) {
        if (answers[student][question] == key[question]) score++;
      }
      System.out.println("Student " + (student + 1) + ": " + score);
    }
  }

}
\end{lstlisting}
\note{Adapted from the supplied ISB campus folder: Multi{-}dimensional arrays.pptx, slides 71, 72. Uses the first three of the ten supplied students; retains the original five{-}question key and answers. Replaced typographic quotes with Java quotes. Complete file: ISB\_Java\_Examples/IsbLecture22.java.}
\end{frame}

\begin{frame}[fragile,t]{Worked trace and expected result}
\small\renewcommand{\arraystretch}{1.4}\rowcolors{2}{pale}{white}
\begin{tabularx}{\textwidth}{>{\raggedright\arraybackslash}X>{\raggedright\arraybackslash}X>{\raggedright\arraybackslash}X}\rowcolor{navy}
\color{white}\bfseries Student & \color{white}\bfseries Correct question numbers & \color{white}\bfseries Score\\
1 & 2, 4, 5 & 3\\
2 & 1, 2, 5 & 3\\
3 & 3, 5 & 2\\
\end{tabularx}\par\vspace{2mm}
\begin{block}{Expected console output}\ttfamily\footnotesize Student 1: 3\par Student 2: 3\par Student 3: 2\end{block}
\note{Adapted from the supplied ISB campus folder: Multi{-}dimensional arrays.pptx, slides 71, 72. Uses the first three of the ten supplied students; retains the original five{-}question key and answers. Replaced typographic quotes with Java quotes. Trace and expected values independently worked for this edition.}
\end{frame}

\begin{frame}[fragile,t]{Extend the ISB example}
\begin{alertblock}{Practice}Which assumption must be checked when rows come from an input file?\end{alertblock}\vspace{3mm}\begin{exampleblock}{Explain your reasoning}Every student row must have the same number of answers as the key; reject malformed row lengths.\end{exampleblock}
\note{Adapted from the supplied ISB campus folder: Multi{-}dimensional arrays.pptx, slides 71, 72. Uses the first three of the ten supplied students; retains the original five{-}question key and answers. Replaced typographic quotes with Java quotes. Answer: Every student row must have the same number of answers as the key; reject malformed row lengths.}
\end{frame}
\begin{frame}[fragile,t]{Lecture summary}
\begin{columns}[T,onlytextwidth]\begin{column}{0.60\textwidth}
\begin{itemize}
\item and access a table of values
\item Traverse rows and columns
\item row and column totals
\end{itemize}
\end{column}\begin{column}{0.35\textwidth}\small
\begin{itemize}
\item Traverse every element with nested loops to accumulate the grand total.
\item For this square matrix, add data[i][i] in a separate loop.
\item Report each result with a distinct label.
\end{itemize}
\end{column}\end{columns}
\note{Ask students to give one concrete example for the concepts listed. Use the worked solution to connect the topics.\par Syllabus reading: Deitel Ch 7. CLO-2.}
\end{frame}

\begin{frame}[fragile,t]{After-class practice}
\begin{columns}[T,onlytextwidth]\begin{column}{0.60\textwidth}
Represent marks for three students in two quizzes. Print each student total and each quiz average.
\end{column}\begin{column}{0.35\textwidth}\small
Syllabus reading\par Deitel Ch 7

Next lecture: Ragged and multidimensional arrays
\end{column}\end{columns}
\note{Reading labels follow the supplied outline. Check topic headings against the available textbook edition. Require a program and expected results for the practice task.\par Syllabus reading: Deitel Ch 7. CLO-2.}
\end{frame}
\end{document}
